% Pursuing Stacks (cotes 134-2 to 134-8): archive cote 134-5, site page 66 (PDF page 67), typescript page 253.
% typescript. Draft compared: ps4.tex:3321-3397.
% AI-transcribed by Claude Opus 5.5 (claude-opus-5-5, adaptive thinking, effort high), Message Batches API, batch ps2 (msgbatch_01NVwr1Yuo8qJf5rykeGBXPK), 2026-11-12; not yet checked by eye.
\scanpage{134-5/66}
\mspage{253}

it is ns that A be aspheric,  and that the equivalent conditions of th.1
be satisfied.

Indeed,  if follows at once from the definitions \xout{that} and from the
fact that the asphericity structure $(M,M_{as})$ is aspheric i.e. \inserted{that} $e_M$ is \xout{\illegible{1}}
aspheric,  that  i is totally $M_{as}$-aspheric iff it is locally $M_{as}$-aspheric
(i.e. condition (ii) of theorem 1),  and \inserted{if} moreover A is aspheric,  hence
the corollary.

\uline{Corollary} 2.  In order for i to be a test functor,  it is ns that A be
a test category,  that $(M,M_{as})$  be modelizing,  and that the equivalent
conditions of theorem 1 be satisfied.

This follows from cor.1  and the reformulation of the notion of a
test-functor,  given \xout{in xth} p. 245.

This corollary contains \xout{that} the main result of section 65 as a parti-%
cular case,  when taking M = (Cat)  with the usual contractibility struc-%
ture,  \xout{except that} and $C = \hand{\{\Delta}_1\hand{\}}$ ,  except that in loc. cit.  we did not
have to assume beforehand that A be a test category,  but only that A is
aspheric:  this condition,  plus condition (iii) above (namely,  $i^{\hand{*}}(\hand{\Delta}_1)$
locally aspheric)  implies already that A is a test category.  In order
to get also this extra result,  we state still another corollary:\hand{(*)}:

\uline{Corollary} 3:  Assume
\[ \underline{I} = (I, \hand{\delta}_0, \hand{\delta}_1\hand{,} \hand{\mu}) \quad , \text{ with } I \hand{\unsure{\in}} \mathrm{Ob}\, C \ , \]
\margin{multiplicati- ty wholly unne- \unsure{cessary} \unsure{and} \unsure{enough that} $I \in M_c$ cf. p. \unsure{348} \unsure{remark 1}}
is a \uline{multiplicative interval} in M ,  i.e. an interval endo\xout{\unsure{w}}wed with a
multiplication $\hand{\mu}$, admitting $\hand{\delta}_0$ as a left unit and $\hand{\delta}_1$ as a left zero ele-%
ment (cf. section 49, p.108  and section 51, p.120  - where such an inter-%
val was provisionally called a ``contractor'').  Assume moreover that for
any x in $M_c$ ,  the two compositions  $x \hand{\to} e_M \hand{\overset{\delta_0,\delta_1}{\rightrightarrows}} I$  are distinct\hand{(} \xout{for} which
is the case for instance if  $\mathrm{Ker}(\hand{\delta}_0,\hand{\delta}_1)$  exists in M and is a strict ini-%
tial object $\emptyset_M$  of M  (i.e. an initial object such that any map $x \hand{\to} \emptyset_M$ in
M is an isomorphism),  and \inserted{moreover} $\emptyset_M \hand{\notin} M_c$ \hand{)}. Then the conditions of th.1  imply
that A is a local test category,  and hence a test category provided A is
aspheric.

Indeed,  $i^*(\underline{I})$  is a multiplicative interval in $A^\wedge$ which is locally
aspheric,  and (as follows immediately from the assumptions on $\underline{I}$) separa-%
ting  -  hence $A^\wedge$ is a local test category.

\margin{66}
